\documentclass[11pt]{scrartcl}

\usepackage[stixtwo]{fontsetup}
\usepackage{amsmath}

\begin{document}
$a_{2N} = 0$. Indoklás: az origó középpontú egységkörbe írt
szabályos $2N$-szög csúcsaiba mutató vektorok összege az elrendezés
szimmetriája miatt a nullvektor, és $2N \cdot a_{2N}$ éppen ezen vektorok
vízszintes koordinátáinak összege.

Hasonlóan igazolható, hogy ha $n$ a $2N$ többszöröse, akkor $a_n = 0$.

Így ha $n = \alpha \cdot 2N + \beta$ ($\alpha, \beta \in \mathbb{N}$,
$\beta < 2N$), akkor
\[
	a_n
	= \frac{1}{n} \sum_{k = 1}^{n} \cos\left(\frac{k\pi}{N}\right)
	= \frac{1}{n} \sum_{k = 1}^{\beta} \cos\left(\frac{k\pi}{N}\right),
	\qquad \text{ezért} \qquad
	|a_n| \le \frac{2N}{n} \to 0.
\]
Tehát az $(a_n)$ sorozat konvergens, határértéke $0$, így
egyetlen torlódási pontja a $0$.

Az $a_n$ képletében szereplő összegzés el is végezhető, ha észrevesszük, hogy
$\cos \frac{k\pi}{N} = \operatorname{Re}(q^k)$, ahol $q = e^{i \frac{\pi}{N}}$,
tehát $n a_n$ egy $q$ hányadosú, $n$ tagú mértani sor összegének valós része,
így
\(
a_n = \frac{1}{n} \operatorname{Re}\left(q \frac{1 - q^n}{1 - q}\right)
\).
Innen is megkapható, hogy
\(
|a_n| \le \frac{1}{n} \left| q \frac{1 - q^n}{1 - q} \right|
\le \frac{2}{n |1 - q|} \to 0
\).
\end{document}
